% Felipe Muñoz Casasola
% This is free and unencumbered software released into the public domain.

% Anyone is free to copy, modify, publish, use, compile, sell, or
% distribute this software, either in source code form or as a compiled
% binary, for any purpose, commercial or non-commercial, and by any
% means.

% In jurisdictions that recognize copyright laws, the author or authors
% of this software dedicate any and all copyright interest in the
% software to the public domain. We make this dedication for the benefit
% of the public at large and to the detriment of our heirs and
% successors. We intend this dedication to be an overt act of
% relinquishment in perpetuity of all present and future rights to this
% software under copyright law.

% THE SOFTWARE IS PROVIDED "AS IS", WITHOUT WARRANTY OF ANY KIND,
% EXPRESS OR IMPLIED, INCLUDING BUT NOT LIMITED TO THE WARRANTIES OF
% MERCHANTABILITY, FITNESS FOR A PARTICULAR PURPOSE AND NONINFRINGEMENT.
% IN NO EVENT SHALL THE AUTHORS BE LIABLE FOR ANY CLAIM, DAMAGES OR
% OTHER LIABILITY, WHETHER IN AN ACTION OF CONTRACT, TORT OR OTHERWISE,
% ARISING FROM, OUT OF OR IN CONNECTION WITH THE SOFTWARE OR THE USE OR
% OTHER DEALINGS IN THE SOFTWARE.

% For more information, please refer to <https://unlicense.org/>

\documentclass{article}

\usepackage[spanish]{babel}
\usepackage[margin=3cm]{geometry}
\usepackage{microtype}
\usepackage[pdfusetitle]{hyperref}

\title{Perfiles de trabajo}
\author{Felipe Muñoz Casasola}
\date{19 de marzo de 2026}

\begin{document}
    \maketitle
    \clearpage
    \tableofcontents
    \clearpage

    \section{Gestión de perfiles}
    En mi carpeta \texttt{$\sim$/Documentos/Gestión\ del\ sistema/Perfiles} se
    encuentran unos guiones en Bash que inician automáticamente el perfiles
    deseado. Cada guion tiene un nombre descriptivo de lo que almacena.
    Debe estar instalado el paquete \texttt{podman-toolbox}, y se puede
    comprobar que ha funcionado porque el nombre de la máquina a partir de
    ese momento es toolbox.

    \section{Instalación}
    Se instaló desde el gestor de paquetes con
    \texttt{sudo apt install podman-toolbox}.

    \section{Creación de un nuevo perfil}
    Para crear un nuevo perfil primero tengo que ejecutar el comando:

    \texttt{toolbox create --distro ubuntu/arch/fedora --release 22.04/rolling/43f}

    Y para meterme en el contenedor recién creado hago:

    \texttt{toolbox enter ubuntu/arch/fedora-toolbox-22.04/rolling/43f}

    Una vez instalados los paquetes necesarios dentro del contenedor y con
    el entorno que desee me puedo salir. Y ejecutar el siguiente comando
    para guardar los cambios:

    \texttt{podman commit --squash ubuntu-toolbox-22.04 localhost/my-ubuntu-toolbox:22.04}

    He sacado esta información (donde estará actualizada) de
    \href{https://containertoolbx.org/doc/}{https://containertoolbx.org/doc/}.
    Y de \href{https://containertoolbx.org/install/}
    {https://containertoolbx.org/install/}. Para más información consultar su
    página web oficial: \href{https://containertoolbx.org}
    {https://containertoolbx.org}.
\end{document}
